% ---------------------------------------------------------------------------
\section{Limitations}
\label{sec:limitations}

\paragraph{No prospective clinician study.}
Our validity anchor is DENTEX's dentist-verified ground truth (each image was annotated by a
final-year dental student and verified by an expert dentist with over 15 years of experience),
together with automated checks of reasoning grounding by a cross-family LLM judge. This
validates against expert labels, not against a prospective, blinded review by practicing
clinicians on new cases. A reader study is the natural next step.

\paragraph{Pending trained-policy results.}
The training-side evaluations (SFT, GRPO, and ablations on the trained agent) are designed,
frozen, and implemented, but the runs are not complete. Those cells are marked as pending
instead of extrapolated. The completed evidence (zero-shot baselines, grounding-tool
validation, and the corpus-level mechanism analysis) stands independently, but end-to-end
claims cannot be made until the runs are complete.

\paragraph{Synthetic perturbations are not real detector errors.}
The tiered perturbation draws offsets from fixed configured ranges, not from the grounding
model's actual error distribution. This keeps demonstrations stable as the detector improves
(\cref{sec:perturbation}), but the training distribution of box errors is synthetic: uniform
random offsets with random signs, unlike the structured errors a real detector makes (for
example adjacent-tooth confusions and root-tip truncations). The evaluation-time
grounding-perturbation probe (\cref{sec:ablations}) measures robustness to a related synthetic
distribution. Robustness to the detector's true error modes is best assessed by swapping in
improved detector weights, which our decoupling makes safe for the trace corpus but which does
not replace measuring it.

\paragraph{Demonstrated vs.\ learned correction.}
\cref{sec:mechanism} measures the \emph{teacher's} behavior in the demonstrations, which is the
skill the student is trained to acquire, not the trained policy's behavior at inference. The
policy-side version of the same analysis is the grounding-perturbation probe, which awaits the
trained checkpoint. Demonstrated behavior is necessary but not sufficient for learned behavior.

\paragraph{Scale and scope.}
The corpus holds 3{,}694 traces over two institutions and four DENTEX diagnosis classes (plus
Tufts' four-class taxonomy). Prior work suggests that a few hundred verified examples can
suffice for medical RL fine-tuning \citep{hu2025medvlmr1}, but the corpus is modest in size and
the DENTEX test set contains only 50 images, so confidence intervals on it are wide; we report
them rather than over-reading point estimates. The tool suite covers radiologist-like
inspection primitives. Measurement and calibration tools (lesion sizing, periodontal probing
analogues) and anomaly detection beyond the four classes are out of scope.

\paragraph{Automated verification is not expert review.}
The verifier is a frontier model from a different family than the generator, audited against
expert labels, and its rejection taxonomy shows that it catches real failures. Automated
verification nevertheless has blind spots that human clinical review would not. We mitigate
this with cross-family verification, ground-truth anchoring of every correctness claim, and
release of the full corpus and audit trail so that others can re-audit it.

\paragraph{Harness and serving evidence.}
The rules-only condition is a prompt-identity control, not a proposed source of diagnostic
improvement. Skills add procedural context and may help or distract, so their effects require
paired dev-split measurements. A BF16 merged target is numerically different from a 4-bit
QLoRA serving target, so speed or equivalence cannot be attributed to speculation across those
targets. Multimodal and MTP support and memory feasibility must be established for the pinned
runtime and hardware; a recognized configuration flag is not evidence of a completed run.
Neither the harness nor speculative decoding establishes clinical validity.
